\documentclass{thecours}
\begin{document}
\bigpart{Distributivité simple}
\proposition{distributivite}
\par\noindent{\bfseries Remarque :}
Cette formule permet de transformer un produit en une somme et vice versa.\vskip10pt
\par\noindent{\bfseries Exemple :}
Développer les expressions (transformer le produit en somme)~:
$$A=5(7+x)\hskip50ptB=8(4-3x)$$
\par\noindent{\bfseries Exemple :}
Factoriser les expressions (transformer la somme en produit)~:
$$A=35x+5y\hskip50ptB=300a-200b$$
\par\noindent{\bfseries Exemple :}
$$25\times(4+4+4+4+4+4+4)$$

\bigpart{Calcul avec les fractions}
% \proposition{calcul_fractionnaire}

\bigpart{Distributivite double}
\proposition{distributivite_double}
\par\noindent{\bfseries Exemple :}
$$A=(3x+1)(2x-1)\hskip50ptB=(\dfrac{3}{2}x-4)^2$$
\newpage
\bigpart{Identitié remarquable}
\proposition{identite_remarquable_3}
\par\noindent{\bfseries Exemple :}
$$A=9x^2-64\hskip50ptB=(x+2)(x-2)$$
\end{document}
